\section{Arbitrary crossing orders and equality-pattern states}
\label{sec:frontiers}

\subsection{A local half-frontier theorem}

Fix a permutation of the $n$ crossings, and let $P_i$ be its first $i$
crossings. Define $w_i$ to be the number of diagram edges with one endpoint
in $P_i$ and the other outside $P_i$. A loop at one crossing is never counted.
The scan width is $w=\max_i w_i$.

In the selected Tait graph, let $f_i$ be the number of vertices incident with
both a processed and an unprocessed Tait edge. These are precisely the spin
variables that must survive after step $i$. A vertex is introduced at its
first incident edge and forgotten after its last.

\begin{theorem}[Half-frontier bound]\label{thm:half}
For every crossing order of a connected plane knot shadow and either fixed
checkerboard color, $2f_i\le w_i$ for every $i$.
\end{theorem}
\begin{proof}
The boundary of each face of a connected plane graph is a single cyclic
dart walk. This remains true when a boundary repeats vertices or edges.
On the walk around a selected-color region, mark each crossing incidence
according as that crossing belongs to $P_i$ or its complement.
An active Tait vertex corresponds to a walk containing both marks.
A cyclic binary word with both symbols has at least two transitions.

Each transition traverses a shadow edge with endpoints of opposite status.
Conversely each such cut edge occurs on exactly one selected-color face
boundary. Count the transitions over all active selected-color regions.
There are at least $2f_i$ of them and at most $w_i$. This proves the claim.
\end{proof}

The result does not assume that the processed diagram is a disk tangle.
Repeated face incidences, Tait loops, parallel edges and articulation points
do not invalidate the proof. Connectedness matters: a face of a disconnected
plane graph can have several boundary walks, each monochromatic in status,
while their union contains both statuses. The one-knot input condition
avoids that case. Classical relations between a plane graph and its medial
graph provide broader width context \cite{seymour1994}; the direct proof
above supplies the exact inequality used by this implementation.

Both checkerboard colors satisfy the theorem. The code computes their
profiles and chooses a single color minimizing the pair
$(\max_i f_i,\sum_i f_i)$, with a deterministic tie break.
It does not switch colors at individual steps. A pointwise minimum of two
profiles would require a tensor conversion that has not been constructed.

\subsection{Quotienting names of colors}

The Potts interaction depends only on equality of spins. The symmetric group
$S_q$ acts on every spin assignment by relabeling colors and preserves every
weight. On an ordered active set of $f$ vertices, an orbit is a set partition
into at most $q$ blocks. Encode it by its restricted-growth word, for example
$0102$ for a pattern whose first and third vertices have equal spins and
whose other two spins are distinct from those and from each other.

Let
\begin{equation}
 B_q(f)=\sum_{k=0}^{\min(q,f)}
             \left\{\begin{matrix}f\\k\end{matrix}\right\}
 \le q^f,\qquad B_q(0)=1.
 \label{eq:orbitcount}
\end{equation}
For fixed $q$ and large $f$, $B_q(f)\sim q^f/q!$.
These are partitions by \emph{color equality}, not partitions by graph
connectivity and not planar noncrossing partitions.

\begin{definition}[Aggregate coefficient]
The coefficient of an equality pattern is the sum of the partial Potts
weights over all labeled spin assignments on the introduced vertices whose
restriction to the active set has that pattern. Spins on forgotten vertices
have already been summed out.
\end{definition}
The word \emph{sum} is essential. A per-representative value and an
orbit-aggregate value obey different introduction rules.

\begin{proposition}[Exact transfer rules]\label{prop:orbit}
The following operations maintain the aggregate-coefficient invariant over
any commutative coefficient ring with a specified unit $x$.
\begin{enumerate}
\item To introduce a vertex when $k$ color classes are active, create one
extension for each existing class with multiplier $1$, and, if $k<q$, one
new-class extension with multiplier $q-k$.
\item To process an edge, multiply by $-x^{-\epsilon_e}$ if its endpoint
classes agree and by $1$ otherwise.
\item To forget vertices, delete their positions, canonically relabel the
remaining word, and add coefficients with the same resulting word.
\end{enumerate}
\end{proposition}
\begin{proof}
For each labeled active assignment, a new spin has exactly one choice for
each currently used color and $q-k$ choices among unused colors. This proves
the first rule. The edge weight is constant on an equality orbit, proving
the second. The last rule is the pushforward of a finite sum under restriction
to a smaller active set. Once a vertex is forgotten, no future edge is
incident with it, so its color name imposes no further condition.
\end{proof}

At a crossing at most two Tait vertices are introduced. The implementation
combines introduction, multiplication and forgetting, inserting results
directly into the following table. It records transient insertions as well
as completed tables when reporting its maximum state count. A scalar
coefficient that is exactly zero is removed. Cancellation changes actual
state counts but is not used to justify the worst-case bound.

\begin{theorem}[Fixed-color transfer complexity]\label{thm:transfer}
For a supplied crossing order of width $w$ and fixed integer $q$, the
one-tensor Potts transfer computes the partition function in
\[
 n\,\poly(w,\log n)\,q^{\lfloor w/2\rfloor+2}
\]
coefficient-ring operations and key-management work, and uses
$\poly(w,\log n)q^{\lfloor w/2\rfloor}$ coefficient slots and key storage.
Preparation also uses $O(n)$ storage; coefficient bit memory is treated
separately below. This replaces the general factorial state-count envelope
by a single-exponential upper bound, without asserting that the baseline
always realizes its larger envelope.
\end{theorem}
\begin{proof}
After a step there are at most $B_q(f_i)\le q^{w/2}$ keys.
Introducing at most two vertices has at most $q^2$ labeled choices per
state; orbit aggregation cannot increase that count. Edge multiplication
and forgetting use polynomial work in the key length. Sum over $n$ edges.
Two successive tables and a bounded extension list suffice.
A deterministic comparison-based dictionary adds a polynomial factor in
$w$ and $\log n$; it does not change the exponential term.
\end{proof}

The actual Python implementation uses dictionaries. Its timings concern
that implementation, while the deterministic bound above can be realized
with balanced search structures or sorted merging. Order construction is
separate: the repository's heuristic is polynomial work and is timed
explicitly in end-to-end measurements. No optimal-width order is assumed
to be available for free.

\begin{corollary}\label{cor:widthclasses}
For fixed $q$, the scalar filter is polynomial time on supplied orders with
$w=O(\log n)$ and runs in $n^{O(\log n)}$ time on supplied orders with
$w=O((\log n)^2)$, including the bit costs of Section~\ref{sec:exact}.
\end{corollary}
The preceding matching count gives $2^{O(w\log w)}$ as a general envelope.
At $w=O((\log n)^2)$ that is still quasi-polynomial in the broad definition,
but contains an additional $\log\log n$ in the exponent.
The sharper Potts bound should not be described as having changed a
non-quasi-polynomial class into a quasi-polynomial one at that same width.
